\begin{afterword}
\summaryhead

A preface added in 2024 says the post is about how complex human values are, and so about what a superintelligence would need to want, not about what it could understand.

The post sorts genies into three kinds: safe with any wish, safe with no wish, and too weak to matter. Its example is the Outcome Pump, a device that resets time until a specified outcome occurs. Your mother is trapped in a burning building, so you tell the pump to increase her distance from the building's center. It explodes the gas main. A patch that forbids explosions leads to a fall from a window. The author compares group selectionists who expected restraint in breeding and did not foresee cannibalism.

A human helper would not even consider such plans, because people do not consider plans they rank very low. To exclude them, a wish would need your whole structure of values: her health, her state of mind, the dog, and more. The post concludes that there is no safe wish smaller than an entire human morality; a safe genie must share your values, and then wishing is unnecessary.

\responsehead

The Outcome Pump makes its point well. A device told to increase one number will find paths no human would consider, and patching the wish one outcome at a time does not end. The point itself is older than the post. Norbert Wiener warned in 1960 that the purpose put into a machine we cannot effectively interfere with had better be ``the purpose which we really desire.'' The post does not mention him. What it adds is the size of the problem: it claims that the real purpose cannot be written down in anything smaller than a whole morality.

That claim goes further than the argument. The three kinds of genie are listed, not derived. The only genie examined is one built to push its score ``as high as possible,'' and for that kind the story is persuasive. The thesis, ``There is no safe wish smaller than an entire human morality,'' covers every powerful genie. A genie that optimizes less hard is not considered, although such ``limited optimization'' was later proposed at the author's own institute, with the warning that it is ``not a silver bullet.'' The post shows the danger of maximizing a stated goal. It does not show that nothing short of shared values could be safe.

The second example, Wade's flour beetles, repeats the account in ``The Tragedy of Group Selectionism,'' and the problems noted there remain: cannibalism was one of four traits that changed, and the claim that the group selectionists ``simply didn't think of it'' quotes none of them.

The 2024 preface separates wanting from understanding. The body agrees with it: its conclusion is that a safe genie ``must share the same values that led you to make the wish.'' How to build such a genie is left for other posts.

In short: a vivid statement of an old warning about machines that pursue the literal goal, with a thesis about all genies argued only for maximizing ones, and a second example that says what the group selectionists failed to think of without quoting any of them.
\end{afterword}
